\documentclass[12pt]{article}
\usepackage[utf8]{inputenc}
\usepackage[T1]{fontenc}
\usepackage{hyperref}
\title{Bilaspuri Language Resources for K-12 Curriculum Development}
\author{Compiled by Claude Sonnet 5.5 (claude-sonnet-5-5) for Andr\'as Kornai}
\date{}
\begin{document}
\maketitle

\section{Introduction}
Bilaspuri (Kahluri; ISO 639-3 kfs) is a Western Pahari language of Bilaspur district, Himachal Pradesh, and nearby parts of Punjab. The 2011 Census counted about 295,805 speakers but classes it under Punjabi. Little material was found.

\section{Primary Educational Resources}
None found in the search results.

\section{Linguistic and Reference Materials}
\begin{itemize}
\item Linguistic Survey of India, Himachal Pradesh volume, with a chapter on Bilaspuri/Kahluri (pp.\ 713--745) covering grammar and sociolinguistic profile: \url{https://ruralindiaonline.org/en/library/resource/linguistic-survey-of-india---himachal-pradesh/} (PDF: \url{https://beta-old.ruralindiaonline.org/v1/uploads/Language_Survey_of_India_Himachal_Pradesh_9a8d9d3690.pdf})
\item Encyclopedic overview: \url{https://en.wikipedia.org/wiki/Bilaspuri}
\end{itemize}

\section{Religious and Community Resources}
None found.

\section{Digital and Audio Resources}
None found.

\section{Curriculum Development Context}
The language is reported to be restricted largely to home communication, and the traditional Takri script is mentioned alongside Devanagari. No school materials were found.
\end{document}